\begin{idea}[title={Откуда этот материал}]
Дальнейшие параграфы этой главы на разобранных лекциях подробно не
записывались — вычислительная схема симплекс-метода была отнесена к
методичкам. Ниже она восстановлена полностью, в обозначениях курса, вместе с
доказательствами; это материал экзаменационного вопроса «Симплекс-метод.
Метод искусственного базиса».
\end{idea}

\subsection{Метод перебора крайних точек}
\label{subsec:enumeration}

Следствие из теоремы~\ref{th:optimal-is-support} даёт прямой, хотя и
неэффективный, способ решения.

\begin{algo}[title={Перебор крайних точек}]
\begin{steps}
  \item Выбрать набор базисных индексов $N_k \subset N$, $\abs{N_k} = m$.
  \item Построить квадратную матрицу $A[M, N_k]$.
  \item Проверить $\det A[M,N_k]$. Если $\det A[M,N_k] = 0$ — набор
        отбрасывается; иначе идти дальше.
  \item Решить систему $A[M,N_k]\,x[N_k] = b[M]$.
  \item Небазисные переменные положить нулю: $x[N \setminus N_k] = 0$.
  \item Проверить $x[N] \ge 0$. Если есть отрицательные компоненты — точка
        отбрасывается; иначе получена допустимая крайняя точка.
  \item Вычислить $F(x) = c\T[N]\,x[N]$.
  \item Повторить шаги 1--7 для всех наборов из $m$ столбцов.
  \item Среди найденных допустимых крайних точек выбрать точку с минимальным
        $F(x)$.
\end{steps}
\end{algo}

\begin{idea}
Метод конечен: наборов не больше $\binom{n}{m}$ (лемма~\ref{lem:count}).
Применим, когда задача записана в канонической форме, $\rank A = m$, $m < n$,
допустимое множество непусто и оптимум конечен. Если ни один допустимый
опорный вектор не найден — допустимое множество пусто. Неограниченность снизу
перебор сам по себе не обнаруживает: для этого нужен симплекс-метод.

Практическая негодность метода видна из оценки: при $n = 50$, $m = 25$ число
наборов $\binom{50}{25} \approx 1{,}26\cdot10^{14}$.
\end{idea}

\subsection{Симплекс-метод: вычислительная схема}
\label{subsec:simplex-scheme}

Симплекс-метод строит последовательность опорных векторов
$x^0[N], x^1[N], \dots$ так, чтобы значение целевой функции не возрастало:
\[
  c\T[N]\,x^{k+1}[N] \ \le\ c\T[N]\,x^{k}[N] .
\]

\subsubsection*{Обозначения}

Пусть на $k$-й итерации известен опорный вектор $x^k[N] \in S$ и его базис
$N_k \subset N$, $\abs{N_k} = m$, $\det A[M,N_k] \ne 0$. Обозначим
\[
  B[N_k, M] = A[M,N_k]^{-1} ,
  \qquad
  L_k = N \setminus N_k \quad \text{— небазисные индексы} ,
\]
при этом $x^k[L_k] = 0$, а базисные компоненты находятся из системы
$A[M,N_k]\,x^k[N_k] = b[M]$.

\subsubsection*{Вектор потенциалов и оценки}

Вектор $y^k[M]$ находится из системы
\[
  \big(y^k\big)\T[M]\, A[M,N_k] = c\T[N_k] ,
  \qquad\text{то есть}\qquad
  A\T[M,N_k]\, y^k[M] = c[N_k] ,
\]
откуда $\big(y^k\big)\T[M] = c\T[N_k]\,B[N_k,M]$. Для каждого небазисного
индекса $j \in L_k$ вычисляется \textbf{оценка}
\[
  d[j] = c[j] - \big(y^k\big)\T[M]\, a_j[M] ,
\]
где $a_j[M]$ — $j$-й столбец матрицы $A[M,N]$.

\begin{idea}
Это ровно та же конструкция, что и потенциалы транспортной задачи
(\S\ref{subsec:potentials-theory} и §7): $y^k$ — двойственные переменные,
$d[j]$ — «во сколько обойдётся» ввод переменной $x[j]$ в базис с учётом
вынужденного пересчёта базисных компонент. Отрицательная оценка означает, что
ввод выгоден.
\end{idea}

\begin{theorem}[критерий оптимальности]\label{thm:simplex-opt}
Если для текущего опорного вектора $x^k[N]$ выполняется $d[j] \ge 0$ для всех
$j \in L_k$, то $x^k[N]$ — оптимальное решение задачи ЛП.
\end{theorem}

\begin{proof}
Пусть $x[N] \in S$ — любое допустимое решение. Оба вектора удовлетворяют
системе, $A[M,N]x[N] = b[M] = A[M,N]x^k[N]$, поэтому
$A[M,N]\big(x[N] - x^k[N]\big) = 0$. Разбивая переменные на базисные и
небазисные,
\begin{gather*}
  A[M,N_k]\big(x[N_k] - x^k[N_k]\big) + A[M,L_k]\,x[L_k] = 0 , \\
  x[N_k] - x^k[N_k] = -B[N_k,M]\,A[M,L_k]\,x[L_k] .
\end{gather*}
Рассмотрим разность значений целевой функции:
\begin{align*}
  c\T[N]x[N] - c\T[N]x^k[N]
  &= c\T[N_k]\big(x[N_k] - x^k[N_k]\big) + c\T[L_k]\,x[L_k] \\
  &= -\,\underbrace{c\T[N_k]B[N_k,M]}_{=\ (y^k)\T[M]}A[M,L_k]\,x[L_k]
     + c\T[L_k]\,x[L_k] \\
  &= \sum_{j \in L_k}
     \underbrace{\Big(c[j] - (y^k)\T[M]\,a_j[M]\Big)}_{=\ d[j]}\, x[j] .
\end{align*}
При $d[j] \ge 0$ и $x[j] \ge 0$ сумма неотрицательна, значит
$c\T[N]x[N] \ge c\T[N]x^k[N]$: вектор $x^k[N]$ оптимален.
\end{proof}

\subsubsection*{Направление перехода}

Если существует индекс $j_k \in L_k$ с $d[j_k] < 0$, текущий опорный вектор не
оптимален, и переменная $x[j_k]$ \textbf{вводится в базис}. Направление
$u^k[N]$ строится так:
\[
  u^k[N_k] = B[N_k,M]\,a_{j_k}[M] ,
  \qquad
  u^k[j_k] = -1 ,
  \qquad
  u^k\big[L_k \setminus \{j_k\}\big] = 0 ,
\]
и рассматривается $x_\theta[N] = x^k[N] - \theta\,u^k[N]$, $\theta \ge 0$.

\begin{lemma}\label{lem:simplex-dir}
$A[M,N]\,u^k[N] = 0$, и потому $A[M,N]\,x_\theta[N] = b[M]$ при любом $\theta$.
\end{lemma}

\begin{proof}
По построению
\[
  A[M,N]u^k[N] = A[M,N_k]u^k[N_k] + a_{j_k}[M]\,u^k[j_k]
    + A\big[M, L_k \setminus \{j_k\}\big]\,u^k\big[L_k\setminus\{j_k\}\big] .
\]
Подставляя $u^k[N_k] = B[N_k,M]a_{j_k}[M]$, $u^k[j_k] = -1$ и нулевой остаток,
\[
  A[M,N]u^k[N] = \underbrace{A[M,N_k]B[N_k,M]}_{=\ E}\,a_{j_k}[M] - a_{j_k}[M] = 0 .
  \qedhere
\]
\end{proof}

\subsubsection*{Изменение целевой функции}

\[
  F(x_\theta) = c\T[N]x_\theta[N] = F(x^k) - \theta\,c\T[N]u^k[N] .
\]
Вычислим $c\T[N]u^k[N] = c\T[N_k]u^k[N_k] + c[j_k]u^k[j_k]
= c\T[N_k]B[N_k,M]a_{j_k}[M] - c[j_k]$. Так как
$c\T[N_k]B[N_k,M] = (y^k)\T[M]$,
\[
  c\T[N]u^k[N] = (y^k)\T[M]a_{j_k}[M] - c[j_k] = -\,d[j_k]
  \qquad\Longrightarrow\qquad
  F(x_\theta) = F(x^k) + \theta\, d[j_k] .
\]
Так как $d[j_k] < 0$, при $\theta > 0$ получаем $F(x_\theta) < F(x^k)$.

\subsubsection*{Выбор шага}

Для сохранения $x_\theta[N] \ge 0$ нужно $x^k[i] - \theta u^k[i] \ge 0$ при всех
$i \in N_k$.
\begin{itemize}[nosep,left=0pt]
  \item Если $u^k[i] \le 0$ для всех $i \in N_k$, то неотрицательность не
        нарушается ни при каком $\theta \ge 0$, а $F(x_\theta) \to -\infty$ при
        $\theta \to +\infty$ — \textbf{задача не ограничена снизу}.
  \item Иначе шаг выбирается по правилу минимального отношения
        \[
          \theta_k = \min_{\substack{i \in N_k \\ u^k[i] > 0}}
                     \frac{x^k[i]}{u^k[i]} .
        \]
\end{itemize}
Пусть минимум достигается на $i_k$. Тогда $x_{\theta_k}[i_k] = 0$: переменная
$x[i_k]$ \textbf{выходит из базиса}, $x[j_k]$ входит,
\[
  N_{k+1} = \big(N_k \setminus \{i_k\}\big) \cup \{j_k\} ,
  \qquad
  x^{k+1}[N] = x^k[N] - \theta_k u^k[N] .
\]

\begin{algo}[title={Симплекс-метод}]
Известен начальный опорный вектор $x^0[N]$ и его базис $N_0$.
\begin{steps}
  \item На $k$-й итерации известны $x^k[N]$, $N_k$, $L_k = N \setminus N_k$.
  \item Построить $B[N_k,M] = A[M,N_k]^{-1}$.
  \item Найти $y^k[M]$ из $\big(y^k\big)\T[M]A[M,N_k] = c\T[N_k]$.
  \item Для каждого $j \in L_k$ вычислить
        $d[j] = c[j] - \big(y^k\big)\T[M]a_j[M]$.
  \item Если $d[j] \ge 0$ для всех $j \in L_k$ — $x^k[N]$ оптимален, конец.
  \item Иначе выбрать $j_k \in L_k$ с $d[j_k] < 0$.
  \item Построить $u^k[N_k] = B[N_k,M]a_{j_k}[M]$, $u^k[j_k] = -1$,
        $u^k[L_k\setminus\{j_k\}] = 0$.
  \item Если $u^k[i] \le 0$ при всех $i \in N_k$ — задача не ограничена снизу,
        конец.
  \item Иначе найти
        $\theta_k = \min\limits_{i \in N_k,\, u^k[i]>0} x^k[i]/u^k[i]$;
        пусть минимум достигается на $i_k$.
  \item $x^{k+1}[N] = x^k[N] - \theta_k u^k[N]$.
  \item $N_{k+1} = \big(N_k \setminus \{i_k\}\big) \cup \{j_k\}$.
  \item $k := k+1$, перейти к шагу~2.
\end{steps}
\end{algo}

\subsubsection*{Вырожденность и зацикливание}

Опорный вектор вырожден, если $\abs{\Nzp} < m$ (определение~\ref{def:degenerate}).
Тогда некоторые базисные компоненты равны нулю, и при выборе шага может
получиться $\theta_k = 0$: опорный вектор не меняется, меняется лишь базис. В
вырожденных задачах возможно \textbf{зацикливание} — возврат к уже
встречавшемуся базису. Средство — \textbf{правило Бленда}: при нескольких
кандидатах (и при выборе вводимой, и при выборе выводимой переменной) брать
переменную с наименьшим индексом.

\begin{idea}
Правило Бленда гарантирует конечность: базисов конечное число
(лемма~\ref{lem:count}), а при правиле наименьшего индекса ни один базис не
повторяется. Платой является возможное замедление — обычно выгоднее брать
$j_k$ с наиболее отрицательной оценкой $d[j_k]$, и переключаться на правило
Бленда только при подозрении на зацикливание.
\end{idea}

\subsection{Метод искусственного базиса: обоснование}
\label{subsec:artificial-basis}

Вспомогательная задача построена в~\S4:
\[
  \begin{cases}
    \displaystyle\sum_{i \in M} y[i] \to \min , \\[6pt]
    A[M,N]\,x[N] + E[M,M]\,y[M] = b[M] , \\
    x[N] \ge 0 , \quad y[M] \ge 0 ,
  \end{cases}
  \qquad\text{старт: } x[N] = 0,\ y[M] = b[M] \ge 0 .
\]

\begin{theorem}\label{thm:artificial-basis}
Пусть вспомогательная задача решена и $\big(x^*[N], y^*[M]\big)$ — её
оптимальное решение. Тогда
\begin{enumerate}[label=\arabic*),left=0pt,nosep]
  \item если $\sum_{i \in M} y^*[i] > 0$, то исходная задача не имеет
        допустимых решений;
  \item если $\sum_{i \in M} y^*[i] = 0$, то $x^*[N]$ — допустимое решение
        исходной задачи.
\end{enumerate}
\end{theorem}

\begin{proof}
Для любой допустимой точки вспомогательной задачи $y[M] \ge 0$, поэтому
$\sum_{i\in M} y[i] \ge 0$: минимум не может быть отрицательным.

\medskip
\textbf{Случай 1.} Пусть $\sum_{i\in M} y^*[i] > 0$. От противного: пусть
исходная задача имеет допустимое решение $\hat{x}[N]$, то есть
$A[M,N]\hat{x}[N] = b[M]$, $\hat{x}[N] \ge 0$. Положим $\hat{y}[M] = 0$; тогда
\[
  A[M,N]\hat{x}[N] + E[M,M]\hat{y}[M] = A[M,N]\hat{x}[N] = b[M] ,
\]
то есть пара $\big(\hat{x}[N], \hat{y}[M]\big)$ допустима во вспомогательной
задаче, и значение её целевой функции равно нулю. Это противоречит
положительности минимума. Значит исходная задача несовместна.

\medskip
\textbf{Случай 2.} Пусть $\sum_{i\in M} y^*[i] = 0$. Так как $y^*[i] \ge 0$,
из равенства суммы нулю следует $y^*[M] = 0$. Подставляя в ограничения,
$A[M,N]x^*[N] = b[M]$, а из допустимости $x^*[N] \ge 0$. Значит
$x^*[N] \in S$.
\end{proof}

\begin{algo}[title={Метод искусственного базиса}]
Дана задача $c\T[N]x[N] \to \min$, $A[M,N]x[N] = b[M]$, $x[N] \ge 0$.
\begin{steps}
  \item Добиться $b[M] \ge 0$: строку с $b[i] < 0$ умножить на $-1$.
  \item Ввести искусственные переменные $y[M] \ge 0$.
  \item Составить вспомогательную задачу $\sum_{i \in M} y[i] \to \min$ при
        $A[M,N]x[N] + E[M,M]y[M] = b[M]$, $x[N] \ge 0$, $y[M] \ge 0$.
  \item Взять начальный искусственный базис $x[N] = 0$, $y[M] = b[M]$.
  \item Решить вспомогательную задачу симплекс-методом.
  \item Если $\sum_{i\in M} y^*[i] > 0$ — исходная задача несовместна, конец.
  \item Если $\sum_{i\in M} y^*[i] = 0$, то $y^*[M] = 0$, а $x^*[N]$ — допустимый
        вектор исходной задачи.
  \item Удалить искусственные переменные. Если какие-то из них остались в
        базисе с нулевым значением, заменить их столбцами исходной матрицы
        $A[M,N]$ так, чтобы базисные столбцы остались линейно независимыми.
  \item Полученный базис использовать как начальный для решения исходной задачи
        обычным симплекс-методом.
\end{steps}
\end{algo}

\begin{idea}
Шаг~8 — тонкое место. Искусственная переменная может остаться в базисе с
нулевым значением, если исходная система содержит линейно зависимую строку или
если допустимое множество вырождено. Замена столбца делается так: для
искусственного индекса $i$ ищут небазисный исходный столбец $a_j$ с ненулевым
элементом в $i$-й позиции преобразованной матрицы и вводят его в базис с
$\theta = 0$. Если такого столбца нет, соответствующее уравнение является
следствием остальных и его можно вычеркнуть.
\end{idea}
